% Cato Maior de Senectute (cato-maior-de-senectute) --- English translation
% Translated by Claude (Anthropic) from the Latin (Perseus).
% Style guide: STYLE.md.

\ciceroSection{1} Titus, if I can give you any help, if I can lighten the worry that even now wears at you, lodged and turning in your heart, what shall my reward be? For I may address you, Atticus, in the very lines with which that man — a man of small fortune but great in loyalty — addresses Flamininus, though I know for certain that you are not troubled, as Flamininus was, night and day. For I know the balance of your mind and its evenness, and I understand that what you brought home from Athens was not the surname alone but humanity and good sense as well. And yet I suspect that you too are sometimes stirred, rather deeply, by the very things that move me — but their consolation is a larger matter, and one to be put off to another time. For now I have decided to write you something on old age.

\ciceroSection{2} For I want both you and myself to be relieved of this burden we share, the old age that already presses upon us, or at any rate is coming on; though I know for certain that you, like one who handles everything with measure and wisdom, both bear it and will bear it. But when I wanted to write something on old age, you came to mind as worthy of a gift the two of us might share in common. To me, at least, the making of this book has been so pleasant that it has not only wiped away every annoyance of old age but has even made old age soft and pleasant. Philosophy, then, can never be praised as fully as it deserves: whoever obeys it can pass every stage of life without distress.

\ciceroSection{3} But about the rest I have said much, and will say much again; this book I have sent to you, on old age. I have assigned the whole discourse not to Tithonus, as Aristo of Ceos did — there would be too little authority in a myth — but to old Marcus Cato, so that the argument might carry greater weight; and in his presence I make Laelius and Scipio marvel that he bears old age so easily, with Cato answering them. If he should seem to argue more learnedly than he usually did in his own books, set it down to Greek literature, of which it is well known he became a passionate student in his old age. But why say more? Cato's own words will now set out everything I think about old age.

\ciceroSection{4} "Time and again," said Scipio, "Gaius Laelius here and I marvel at your wisdom, Marcus Cato — at how outstanding and complete it is in all things, but above all at this: that I have never sensed old age weighing on you, when to most old men it is so hateful that they say they carry a load heavier than Etna." It is no very hard thing, said Cato, that you two seem to marvel at, Scipio and Laelius. For those who have no resource in themselves for living well and happily, every age is a burden; but those who seek all their goods from within themselves can find nothing bad in what nature's necessity brings. Old age belongs first of all to this class. Everyone hopes to reach it, and then complains once it is reached — such is the inconstancy and perversity of folly. They say it creeps up sooner than they had expected. First, who forced them into that false expectation? For how does old age creep up on youth any sooner than youth creeps up on boyhood? Then how would old age weigh less on them if they were living their eight-hundredth year rather than their eightieth? For however long the years that have passed, once they had slipped away, no consolation could soothe a foolish old age.

\ciceroSection{5} And so, if you are in the habit of marveling at my wisdom — and I only wish it were worthy of your good opinion and of my surname — our wisdom lies in this: that we follow nature as the best of guides, as though she were a god, and obey her. It is not likely that she, having laid out the other parts of life so well, has neglected the final act like some lazy poet. Yet there had to be some final stage, and, as with the berries of trees and the fruits of the earth, something withered and ready to fall in its ripe season — and this the wise man must bear gently. For what else is making war on the gods, in the manner of the Giants, but fighting against nature?

\ciceroSection{6} "And yet, Cato," said Laelius, "you would do us a great kindness — I can promise it for Scipio too — if, since we hope, and certainly wish, to grow old, we could learn from you well in advance the surest methods for bearing the weight of advancing years." I will indeed, Laelius, said Cato, especially since, as you say, it will please you both. "We do truly wish it," said Laelius, "and, if it is no trouble, Cato, since you have completed a kind of long road that we too must travel, we should like to see what the place you have reached is like."

\ciceroSection{7} I will do as I can, Laelius. For I have often sat through the complaints of men my own age — and like flocks with like, as the old proverb has it. Men of consular rank, Gaius Salinator and Spurius Albinus, more or less my contemporaries, used to lament now that they were cut off from the pleasures without which they thought life was nothing, now that they were scorned by those who had once cultivated them. To me they seemed to be blaming the wrong target. For if this happened through the fault of old age, the same things would befall me and every other older man — yet I have known the old age of many men who passed it without complaint, who were not at all sorry to be loosed from the chains of appetite, and were not despised by their own people. No, the fault in all complaints of this kind lies in character, not in age. The old who are temperate, neither difficult nor unkind, pass a tolerable old age; harshness and unkindness are a trial at every age.

\ciceroSection{8} "It is as you say, Cato," said Laelius, "but perhaps someone might object that old age seems more tolerable to you because of your wealth, your resources, and your standing, and that this cannot fall to many." There is something in that, Laelius, said Cato, but it is by no means everything. There is the story of how Themistocles, in a quarrel, answered a certain man of Seriphus who had said that he owed his brilliance not to himself but to the glory of his country: "By Hercules," he said, "if I were a man of Seriphus, I should never have been famous — and neither would you, if you were an Athenian." The same can be said of old age. For not even to the wise man can old age be light in utter poverty, nor to the fool can it be anything but a burden even in the midst of plenty.

\ciceroSection{9} The most fitting weapons of old age, Scipio and Laelius, are without doubt the practice and exercise of the virtues, which, when cultivated at every age, bring forth wonderful fruits after a long and full life — not only because they never desert you, not even at the very end, though that itself is the greatest thing, but also because the awareness of a life well lived, and the memory of many good deeds, is the sweetest thing there is.

\ciceroSection{10} As a young man I loved Quintus Maximus — the one who recovered Tarentum — old as he was, as if he were my equal in years. For in that man there was a dignity seasoned with courtesy, and old age had not changed his character. To be sure, I began to honor him when he was not yet greatly advanced, though already on in years. For he had been consul for the first time the year after I was born, and as a very young soldier I marched with him to Capua in his fourth consulship, and five years later to Tarentum. Then four years afterward I was made quaestor, an office I held under the consuls Tuditanus and Cethegus, when he, already quite old, spoke in support of the Cincian law on gifts and fees. He waged wars like a young man, though he was plainly well on in years, and by his patience he wore down Hannibal in his youthful exuberance. About him our friend Ennius writes splendidly: "One man, by delaying, restored our cause; he set no rumors above our safety; therefore the more, and more brightly, the man's glory shines now."

\ciceroSection{11} And Tarentum — with what vigilance, with what judgment he recovered it! In my own hearing, when Salinator, who had lost the town and fled to the citadel, boasted to him and said, "It was by my doing, Quintus Fabius, that you recovered Tarentum," he answered with a laugh: "Quite so — for if you had not lost it, I should never have recovered it." Nor was he greater in arms than in the toga. In his second consulship, while his colleague Spurius Carvilius held back, he resisted the tribune of the plebs Gaius Flaminius as far as he could, when Flaminius was dividing up the Picene and Gallic land man by man against the Senate's authority. And as an augur he dared to say that whatever was done for the safety of the republic was done under the best auspices, and whatever was carried against the republic was carried against the auspices.

\ciceroSection{12} I came to know many splendid things in that man, but nothing more admirable than the way he bore the death of his son, a distinguished man who had held the consulship. His eulogy is in our hands, and when we read it, what philosopher do we not look down upon? Yet he was great not only in the public light, before the eyes of his fellow citizens, but more outstanding still within doors, at home. What conversation, what counsels! What a knowledge of antiquity, what mastery of augural law! He had wide learning too, for a Roman: he held in memory every war, not only our own but foreign ones as well. I drank in his conversation so eagerly even then, as though I already divined what in fact came to pass — that once he was gone there would be no one left to learn from.

\ciceroSection{13} Why, then, all this about Maximus? Because surely you see how impious it would be to call such an old age wretched. And yet not everyone can be a Scipio or a Maximus, recalling the storming of cities, battles on land and sea, wars he himself waged, triumphs. There is also a calm and gentle old age that crowns a life lived quietly, purely, and with grace — such as we are told was Plato's, who died at eighty-one, pen in hand; such as Isocrates', who says he wrote the book entitled \textit{Panathenaicus} in his ninety-fourth year and lived five more after it; and his teacher, Gorgias of Leontini, completed a hundred and seven years and never let up in his pursuit and his work. When he was asked why he was willing to remain alive so long, he said, "I have no charge to bring against old age."

\ciceroSection{14} A splendid answer, and worthy of a learned man. For it is fools who lay their own faults and their own blame upon old age — which is not what Ennius did, the man I mentioned just now: "Like a brave horse that often at the last lap won at Olympia, now worn out with age, he rests." He compares his own old age to that of a brave and victorious horse. You can well remember him: it was in the nineteenth year after his death that the present consuls, Titus Flamininus and Manius Acilius, took office, and he himself died when Caepio and Philippus — the latter for the second time — were consuls, the year I, sixty-five years old, argued for the Voconian Law in a strong voice and with sound lungs. At seventy — for that is how long Ennius lived — he bore the two burdens reckoned heaviest, poverty and old age, in such a way that he seemed almost to take delight in them.

\ciceroSection{15} And in fact, when I turn it over in my mind, I find four reasons why old age is thought wretched: first, that it draws us away from active affairs; second, that it makes the body weaker; third, that it strips us of nearly all pleasures; fourth, that it stands not far from death. Let us look, if you please, at how much weight there is in each of these reasons, and how just each is. Old age withdraws us from active affairs. From which? From those carried on by youth and strength? Are there then no tasks proper to the old, to be managed by mind and judgment even when the body is weak? So Quintus Maximus did nothing, then? Nothing, Lucius Paulus, your father, the father-in-law of that excellent man my son? And the rest of the old men — the Fabricii, the Curii, the Coruncanii — when they defended the republic by counsel and authority, were they doing nothing? To the old age of Appius Claudius there was added even this, that he was blind;

\ciceroSection{16} yet when the Senate's opinion was inclining toward peace with Pyrrhus and the making of a treaty, he did not hesitate to say the very words Ennius set down in verse: "Where have your minds, which once stood upright before, turned aside, deranged, from their path?" — and the rest, with the utmost gravity, for the poem is known to you, and besides, Appius's own speech survives. And he did this seventeen years after his second consulship, with ten years having fallen between his two consulships, and having held the censorship before the earlier consulship — from which it is clear that at the time of the war with Pyrrhus he was a very old man. And so it has been handed down to us from our fathers.

\ciceroSection{17} They bring nothing to the case, then, who deny that old age has a part in the conduct of affairs; they are like men who would say the helmsman does nothing in sailing, since some climb the masts, some run along the gangways, some bail out the bilge, while he sits quietly in the stern holding the tiller. He does not do what the young men do; but in fact he does things far greater and better. Great things are accomplished not by strength or speed or quickness of body, but by counsel, authority, and judgment — qualities of which old age is not stripped, but is usually enriched.

\ciceroSection{18} Unless, perhaps, I — who have served as soldier and tribune and legate and consul through every kind of war — seem to you to be idle now, because I wage no wars. And yet I prescribe to the Senate what is to be done, and how. To Carthage, long now plotting mischief, I declare war well in advance; and I shall not cease to fear that city until I have learned she is razed to the ground. May the immortal gods reserve that palm,

\ciceroSection{19} Scipio, for you, that you may finish the work your grandfather left undone — from whose death this is now the thirty-third year, though all the years to follow will take up the memory of that man. He died the year before I was censor, nine years after my consulship, having been elected consul a second time during my own consulship. Tell me, then: if he had lived to his hundredth year, would he have regretted his old age? For he would not have been employing the charge or the leap, nor spears at a distance or swords at close quarters, but counsel, reason, and judgment — and if these were not present in old men, our ancestors would never have called their highest council the Senate. Among the Spartans, indeed, those

\ciceroSection{20} who hold the highest magistracy are, as they are in fact, so also in name called Elders. But if you are willing to read or hear of foreign affairs, you will find that the greatest commonwealths were shaken by young men and sustained and restored by the old. "Tell me — how did you lose your republic, so great, so fast?" For so they put the question in the poet Naevius's play \textit{The Wolf}. Several answers are given, and this above all: "New orators kept coming forward, foolish young upstarts." Recklessness, you see, belongs to the flowering of youth, prudence to its waning.

\ciceroSection{21} But memory weakens. I grant it — unless you exercise it, or unless you are by nature on the slower side. Themistocles had learned the names of all his fellow citizens; do you suppose, then, that as he advanced in years he ever greeted Aristides as Lysimachus? For my part, I know not only the men who are alive, but their fathers too and their grandfathers; nor in reading their epitaphs do I fear, as the saying goes, that I shall lose my memory — for by reading these very stones I return into the memory of the dead. And in truth I have never heard of an old man who forgot where he had buried his treasure. They remember everything they care about: the bail dates fixed, who owes them, and whom they themselves owe.

\ciceroSection{22} What of jurists, what of pontiffs, what of augurs, what of philosophers in old age? How much they remember! The powers of mind remain in old men, provided enthusiasm and industry remain as well — and not in distinguished and honored men only, but in a quiet and private life too. Sophocles composed tragedies into extreme old age; and when, on account of this pursuit, he seemed to be neglecting his household affairs, he was summoned to court by his sons, so that — just as by our custom it is usual to bar fathers who manage their estate badly from control of their property — the jurors might remove him from his estate as if he were senile. Then the old man is said to have recited to the jurors the play he had in hand and had most recently written, \textit{Oedipus at Colonus}, and to have asked whether that poem seemed the work of a man losing his wits;

\ciceroSection{23} When this was read aloud, he was acquitted by the votes of the jurors. Did old age, then, force this man to fall silent in his pursuits? Did it silence Homer, Hesiod, Simonides, Stesichorus, or those I mentioned earlier, Isocrates and Gorgias, or the leaders of philosophy, Pythagoras and Democritus, or Plato and Xenocrates, or, later, Zeno and Cleanthes, or the man whom you yourselves even saw at Rome, Diogenes the Stoic? Was it not rather that in all of them the practice of their pursuits ran on as long as life itself?

\ciceroSection{24} Come, let us set aside these godlike pursuits. I can name Roman farmers from the Sabine country, my neighbors and friends, in whose absence almost no major work on the land ever gets done at all, neither the sowing, nor the gathering, nor the storing of the crops. And yet in their case this is less remarkable, since no man is so old that he does not think he can live another year. But these same men labor at things they know have no bearing on themselves at all: he plants trees to benefit another age, as our own Statius says in his \textit{Synephebi}.

\ciceroSection{25} Nor indeed does the farmer hesitate, however old he is, to answer the man who asks for whom he plants: "For the immortal gods, who willed that I should not only receive these things from my forebears but also pass them on to those who come after." And Caecilius spoke better of the old man who looks ahead to another age than when he wrote this: "By Pollux, old age, if you brought no other fault along with you when you came, this one would be enough: that by living long a man sees many things he does not want to see." And many, perhaps, that he does want; besides, youth too often runs into the very things it does not want. Worse still is this other line of the same Caecilius: "And this I count the most wretched thing in old age, to feel that at this age one has become a burden to others." A delight, rather, than a burden!

\ciceroSection{26} For just as wise old men take delight in young men of good character, and the old age of those who are cultivated and loved by the young is made lighter, so the young rejoice in the precepts of the old, by which they are led toward the pursuit of the virtues; and I am no less aware that I am a delight to you than that you are to me. But you see how old age is not only not languid and idle, but is even busy, always doing something and setting something in motion—of the kind, that is, that each man's pursuit was in his earlier life. And what of those who actually go on learning something new? As we see Solon boasting in his verses, when he says that he grows old learning something fresh each day. And I have done the same, learning Greek letters as an old man; indeed I seized upon them as eagerly as if I were longing to slake a long thirst, so that the very examples you now see me using might be known to me. And when I heard that Socrates had done this with the lyre, I should have liked to do that too—the ancients did learn the lyre—but in letters at any rate I worked hard.

\ciceroSection{27} Nor do I now miss a young man's strength—for this was the second heading, concerning the faults of old age—any more than as a young man I missed the strength of a bull or an elephant. Use what you have, and whatever you do, do it in proportion to your strength. For what utterance could be more contemptible than that of Milo of Croton? When he was already old and watched the athletes training on the track, he is said to have looked at his own arms and to have said, weeping, "But these are dead now." Not so dead as you yourself, you trifler—for it was never from yourself that you won your fame, but from your flanks and your arms. Nothing of the kind from Sextus Aelius, nothing from Tiberius Coruncanius many years before, nothing lately from Publius Crassus, by whom the law was laid down for the citizens, and whose wisdom went on advancing to their last breath.

\ciceroSection{28} As for the orator, I do fear that he may flag in old age, for his function depends not on talent alone but also on lungs and strength. Yet that resonance in the voice somehow even grows brighter in old age—a thing which I myself have not yet lost, and you see my years. Still, a calm and unforced manner of speech becomes an old man, and very often the polished, gentle discourse of an eloquent old man wins itself a hearing of its own accord; and if you cannot deliver it yourself, you can still pass on the art to a Scipio and a Laelius. For what is more delightful than an old age surrounded by the enthusiasm of the young?

\ciceroSection{29} Or shall we not leave old age even the strength to teach the young, to train them, to equip them for every duty of office? And what work can be more splendid than this? For my part, I thought Gnaeus and Publius Scipio, and your two grandfathers, Lucius Aemilius and Publius Africanus, fortunate in their retinue of noble young men; nor should any teachers of the liberal arts be reckoned anything but blessed, however much their strength may have aged and failed. And yet that very failing of strength is more often brought about by the faults of youth than by old age; for a lustful and intemperate youth hands the body over to old age worn out.

\ciceroSection{30} Cyrus, indeed, in Xenophon, in the discourse he gave as he was dying, when he was very old, says that he never felt his old age had become any feebler than his youth had been. As a boy I remember Lucius Metellus, who, having been made pontifex maximus four years after his second consulship, presided over that priesthood for twenty-two years, and was of such sound strength in the last stretch of his life that he did not feel the want of youth. I need say nothing about myself—though that, to be sure, is an old man's privilege, and one granted to our time of life.

\ciceroSection{31} Do you not notice how often, in Homer, Nestor proclaims his own virtues? For he was looking on the third generation of men, and he had no cause to fear that in telling the truth about himself he might seem either too arrogant or too talkative. For indeed, as Homer says, from his tongue flowed speech sweeter than honey, and for that sweetness he needed no strength of body. And yet that famous leader of Greece nowhere wishes to have ten men like Ajax, but ten like Nestor; if that were granted him, he has no doubt that Troy would shortly perish.

\ciceroSection{32} But I return to myself. I am in my eighty-fourth year. I could wish, to be sure, that I could boast the same as Cyrus; but this much I can say: I do not, indeed, have the strength I had as a soldier in the Punic War, or as quaestor in that same war, or as consul in Spain, or four years later, when as military tribune I fought at Thermopylae in the consulship of Manius Glabrio. Yet even so, as you see, old age has not utterly unstrung me, has not laid me low: the Senate house does not miss my strength, nor the Rostra, nor my friends, nor my clients, nor my guests. For I have never assented to that old and much-praised proverb which advises a man to become old early if he wishes to be old for long. I, for my part, would rather be old for a shorter time than be old before I was. And so no one yet has wished to meet with me and found me too busy to see him.

\ciceroSection{33} But I have less strength than either of you. Well, you do not have the strength of the centurion Titus Pontius either: is he, on that account, the better man? Let there only be moderation in the use of one's strength, and let each man strive only so far as he can, and surely he will not be gripped by any great longing for strength. At Olympia Milo is said to have walked the length of the stadium carrying an ox on his shoulders: which, then, would you rather have given to you, this strength of body or the strength of Pythagoras's mind? In short, use this good while it is present; when it is gone, do not long for it—unless, perhaps, young men ought to long for boyhood, and those a little further on in years for their youth. The course of life is fixed, and the path of nature is one and simple, and to each part of life its own fitting season has been given, so that the frailty of children, the fierceness of the young, the gravity of the now settled age, and the ripeness of old age each have something natural about them, something to be enjoyed in its own time.

\ciceroSection{34} You have heard, I imagine, Scipio, what your family's old guest-friend Masinissa does today at the age of ninety: when he has set out on a journey on foot, he never mounts a horse, and when he has set out on horseback, he never dismounts; no rain, no cold can induce him to cover his head; his body is utterly lean and dry, and so he carries out all the duties and offices of a king. Exercise, then, and self-restraint can preserve even in old age something of one's earlier strength. Suppose there is no strength in old age: strength is not even demanded of old age. And so by both law and custom our time of life is exempt from those duties that cannot be discharged without strength. The result is that we are compelled not only to do nothing we cannot do, but not even as much as we can.

\ciceroSection{35} But many old men, you say, are so feeble that they cannot carry out any duty of life, indeed any function at all. Yet that is not a fault peculiar to old age; it is one shared with poor health. How feeble was the son of Publius Africanus, the man who adopted you—what frail health he had, or rather none at all! Had it been otherwise, he would have stood out as a second light of the state, for to his father's greatness of soul he had added a richer learning. What wonder, then, if old men are sometimes infirm, when even the young cannot escape it? Old age must be resisted, Laelius and Scipio, and its faults must be made good by diligence; we must fight against it as against a disease; we must keep an eye on our health,

\ciceroSection{36} take moderate exercise, and consume only as much food and drink as will restore our strength, not crush it. And we must come to the aid not only of the body, but far more of the mind and spirit. For these too, unless you keep dripping oil into them as into a lamp, are snuffed out by old age. And while the body grows heavy with the fatigue of exertion, the mind is lightened by being exercised. For when Caecilius speaks of the foolish old men of comedy, he means the credulous, the forgetful, the slack—faults that belong not to old age as such, but to an idle, listless, drowsy old age. Just as wantonness and lust belong more to the young than to the old—and not to all the young, but to the unprincipled ones—so this senile foolishness, which people are accustomed to call dotage, belongs to the frivolous among the old, not to all of them.

\ciceroSection{37} Four sturdy sons, five daughters, so great a household, so wide a circle of clients—all these Appius governed though he was both blind and old. For he kept his mind taut as a bow, and did not slacken and surrender to old age. He held over his people not merely authority but command: the slaves feared him, the children revered him, all held him dear; in that house the ancestral custom and discipline were alive and strong.

\ciceroSection{38} For old age is honorable on these terms: if it defends itself, if it keeps its own rights, if it is enslaved to no one, if it rules over its own household down to the last breath. For just as I approve a young man who has something of the old in him, so I approve an old man who has something of the young; and whoever follows this principle may be old in body, but in mind he never will be. The seventh book of my \textit{Origines} is in hand; I am collecting all the records of antiquity; I am now hard at work composing the speeches I delivered in the famous cases I defended; I am studying the law of the augurs, the pontiffs, and the civil courts; I make much use of Greek literature, too; and in the manner of the Pythagoreans, for the sake of training my memory, I recall in the evening what I have said, heard, and done each day. These are the exercises of the mind, these the racecourses of the intellect; sweating and laboring at them, I feel no great want of bodily strength. I attend to my friends, I come often to the Senate, and of my own accord I bring before it matters long and deeply pondered, and I defend them by the strength of mind, not of body. And if I were unable to carry these things through, even so my little couch would give me pleasure as I dwelt on the very things I could no longer do. But that I am able to do them is the work of a life so spent. For to one who lives always amid these pursuits and labors, it is never noticed when old age creeps up: thus one's years grow old imperceptibly, without one's feeling it, and life is not suddenly broken off but extinguished by sheer length of time.

\ciceroSection{39} There follows the third charge against old age: that they say it is without pleasures. What a splendid gift of our years, if indeed it takes from us the very thing that is most ruinous in youth! Hear, then, my excellent young men, an old discourse of Archytas of Tarentum, a man among the greatest and most distinguished, which was passed on to me when I was a young man at Tarentum with Quintus Maximus. He used to say that nature had given men no more deadly plague than bodily pleasure, since the appetites greedy for that pleasure were goaded recklessly and without restraint to gratify it.

\ciceroSection{40} From this, he said, spring betrayals of one's country, from this the overthrow of states, from this secret dealings with the enemy; in short, there is no crime, no wicked deed, to whose undertaking the lust for pleasure does not drive a man; and as for debauchery and adultery and every such outrage, they are roused by no other enticements than those of pleasure; and since nature, or some god, had given man nothing more excellent than the mind, nothing was so hostile to this divine gift and bounty as pleasure.

\ciceroSection{41} For where lust holds sway there is no room for self-restraint, and in the kingdom of pleasure virtue can find no foothold at all. To make this easier to grasp, he would bid one imagine someone stirred by as great a bodily pleasure as could possibly be felt; no one, he held, could doubt that as long as he was so taking his delight, he could turn nothing over in his mind, attain nothing by reason, nothing by reflection. Therefore nothing was so detestable and so destructive as pleasure, since, when it was great and prolonged, it snuffed out all the light of the mind. This discourse of Archytas—delivered to the Samnite Gaius Pontius, the father of the man by whom the consuls Spurius Postumius and Titus Veturius were defeated at the battle of the Caudine Forks—our guest-friend Nearchus of Tarentum, who had stayed loyal to his friendship with the Roman people, said he had received from his elders, who added that Plato the Athenian had been present at that conversation; and I find that Plato came to Tarentum in the consulship of Lucius Camillus and Appius Claudius.

\ciceroSection{42} What is the point of all this? To make you understand that, if we could not spurn pleasure by reason and wisdom, we ought to feel great gratitude to old age, which sees to it that we no longer want what we ought not to want. For pleasure obstructs deliberation, is hostile to reason, blinds, so to speak, the eyes of the mind, and has no dealings whatever with virtue. It was against my will that I expelled Lucius Flamininus—brother of that bravest of men, Titus Flamininus—from the Senate seven years after he had been consul; but I judged that his lust had to be branded. For while he was consul in Gaul, he was prevailed upon at a banquet by a courtesan to put to the axe one of those who were in chains, condemned on a capital charge. Under his brother Titus, who was censor just before me, he had slipped through; but to me and to Flaccus a lust so disgraceful and so abandoned could in no way be approved, joining as it did a private outrage to the dishonor of high office.

\ciceroSection{43} I have often heard from my elders, who said that they in turn had heard it as boys from old men, that Gaius Fabricius used to marvel at what he had heard from the Thessalian Cineas, when Fabricius was an envoy at the court of King Pyrrhus: that there was a certain man at Athens who professed himself a sage, and who said that everything we did should be referred to pleasure. And Manius Curius and Tiberius Coruncanius, on hearing this from him, used to wish that the Samnites and Pyrrhus himself could be persuaded of it, so that they might be conquered the more easily once they had given themselves over to pleasures. Manius Curius had lived in the time of Publius Decius, who, five years before Curius was consul, had in his fourth consulship vowed his own life for the republic; Fabricius knew the same Decius, Coruncanius knew him; and these men, judging both from their own lives and from the deed of the Decius I have named, held that there is surely something noble and splendid in nature that is sought for its own sake, and that all the best men pursue when pleasure has been spurned and despised.

\ciceroSection{44} Why, then, so much about pleasure? Because it is not only no reproach to old age, but its highest praise, that it feels no great want of pleasures. It does without lavish feasts and loaded tables and cups passed round; and so it also does without drunkenness and indigestion and sleepless nights. But if some concession must be made to pleasure—since we do not easily withstand its allurements (for Plato, with divine insight, calls pleasure the bait of evil, because men are of course caught by it as fish are)—then, though old age does without immoderate banquets, it can still take delight in modest gatherings. As a boy I often saw Gaius Duellius, son of Marcus, the first man to crush the Carthaginians in a naval battle, coming home from dinner as an old man; he delighted in his wax torch and his flute-player, indulgences which, as a private citizen, he had assumed without precedent: such license did his glory grant him.

\ciceroSection{45} But why speak of others? Let me come back now to myself. To begin with, I have always had my fellow members—the dining-clubs were in fact established in my quaestorship, when the Idaean rites of the Great Mother were brought to Rome—and so I used to dine with them, modestly enough on the whole, though there was a certain heat of youth in it, which fades, year by year, as life advances and everything grows gentler. For I never measured the delight of those gatherings by bodily pleasures so much as by the company of friends and the talk. Our ancestors did well to call a reclining of friends at a meal, since it carries a joining of lives, a \textit{convivium}, a living-together—better than the Greeks, who call this same thing now a drinking-together, now an eating-together, as though they prized most in it what is least worth prizing.

\ciceroSection{46} For my part, on account of the delight of conversation, I take pleasure even in dinners that begin early in the day—and not only with men of my own generation, very few of whom remain, but with yours too, and with you. I owe old age a great debt, which has sharpened my appetite for talk and taken away my appetite for drink and food. And if these very things do still delight someone—for I would not seem to have declared all-out war on pleasure, which perhaps has a certain natural measure—I do not see that even in these pleasures old age lacks all feeling. For my part I delight in the mastership of the feast established by our ancestors, and in the talk that, in the manner of our ancestors, is started by the man at the head of the cup, and in cups that, as in Xenophon's \textit{Symposium}, are small and beaded with dew, and in cooling shade in summer and, by turns, the sun or a winter fire. These very pleasures I am in the habit of pursuing even among my Sabine farms, and every day I fill my table with a gathering of neighbors, which we draw out, with talk of every kind, as far into the night as we can.

\ciceroSection{47} But there is not in the old, you say, so keen a tickling of the pleasures. I believe it—but there is not even a longing for them; and nothing is a torment that you do not miss. Sophocles answered well, when a man asked him, already worn with age, whether he still made use of love: "Heaven forbid!" he said. "I have run from it gladly, as from a savage and raging master." For to those who crave such things, to go without is perhaps hateful and a torment; but to those who have had their fill and are sated, it is more pleasant to go without than to enjoy. And yet he who does not miss a thing does not go without it; and so I say that not to miss it is the more pleasant.

\ciceroSection{48} But suppose the prime of life does enjoy these very pleasures more freely. In the first place, it enjoys, as I have said, small things; and in the second, things that old age, even if it does not have them in abundance, does not lack altogether. As a man delights more in Ambivius Turpio when he watches from the front rows of the theater, yet the man in the back rows takes his delight too, so youth, looking on pleasures from close at hand, perhaps rejoices more, but old age delights in them all the same, watching from a distance—as much as is enough.

\ciceroSection{49} But how great are those other rewards—for the mind to be with itself, its service to lust and ambition, to rivalries and feuds and all the cravings now discharged like a soldier's completed campaigns, and, as they say, to live in its own company! And if it has, besides, some fodder of study and learning, nothing is more delightful than a leisured old age. We saw Gallus, your father's friend, Scipio, all but taking the measure of heaven and earth in his studies. How often the dawn surprised him after he had set to work on some diagram by night, how often night surprised him after he had begun at morning! What delight it gave him to foretell to us, long in advance, the eclipses of the sun and moon!

\ciceroSection{50} And in lighter studies, yet keen ones? How Naevius rejoiced in his \textit{Punic War}, how Plautus in his \textit{Truculentus}, how in his \textit{Pseudolus}! I even saw Livius, an old man, who, having produced a play six years before I was born, in the consulship of Cento and Tuditanus, lived on in years right up to my own youth. Need I speak of the studies of Publius Licinius Crassus, in both pontifical and civil law, or of this Publius Scipio, who within these last few days was made pontifex maximus? And all of these whom I have mentioned we have seen ablaze with these studies as old men. And Marcus Cethegus, whom Ennius rightly called the marrow of Persuasion—with what zeal we saw him put to work in speaking, even as an old man! What pleasures, then, of feasts or games or harlots are to be compared with these pleasures? And these, indeed, are the studies of learning, which for the wise and well-trained grow together with their years—so that what Solon says is honorable, in a line of verse, as I said before, that he grows old learning many things day by day; and surely no pleasure of the mind can be greater than that.

\ciceroSection{51} I come now to the pleasures of the farmer, in which I take incredible delight, which are hindered by no old age and which seem to me to come nearest of all to the life of a wise man. For they have their dealings with the earth, which never refuses its command, and never returns what it has received without interest—now with smaller, but for the most part with larger increase. And yet for me it is not the yield alone, but the very force and nature of the earth herself that delights me. When she has taken the scattered seed into her softened and worked-over lap, first she hides it away, buried—from which harrowing, which does this, takes its name; then, warmed by her own heat and embrace, she opens it out and draws from it a sprouting green, which, propped on the fibers of its roots, grows up little by little, and standing erect on its jointed stalk is now enclosed, as if coming into manhood, in its sheaths; and when it has burst from these, it pours forth fruit ranged in the ordered ear, and against the pecking of small birds is fenced with a palisade of beards.

\ciceroSection{52} Why should I recount the rise, the planting, the growth of the vine? That you may know the rest and solace of my old age, I cannot have my fill of the delight of it. For I pass over the very force of all things that are generated from the earth—which from so tiny a grain of a fig, or from a grape's pip, or from the minutest seeds of the other crops and plants, brings forth such great trunks and branches. The mallet-shoots, the slips, the cuttings, the quicksets, the layers—do they not make anyone delight in wonder? The vine, indeed, which by nature is prone to fall and, unless propped, is carried to the ground, yet, to raise itself, lays hold with its tendrils, as if with hands, on whatever it has found; and as it creeps in winding, wandering descent, the farmer's craft, pruning with the knife, holds it in check, lest it run wild with shoots and spread too far in every direction.

\ciceroSection{53} And so, as spring comes on, in the parts that have been left there appears, at the joints of the shoots as it were, what is called the bud, and from it the rising grape shows itself, which, swelling with the sap of the earth and the heat of the sun, is at first quite sour to the taste, then ripens and grows sweet, and, clothed in its leaves, is neither without a mild warmth nor unprotected from the sun's excessive blaze. What can be more joyful than this for its fruit, or more beautiful for the sight of it? It is not its usefulness alone, as I said before, that delights me, but also its cultivation and its very nature: the rows of props, the yoking of the heads, the binding and layering of the vines, the pruning of some shoots, as I said, and the letting-out of others. Why should I bring forward the irrigations, the trenchings and re-diggings of the field, by which the earth is made far more fertile? Why speak of the usefulness of manuring? I have written of it in the book I composed on country matters.

\ciceroSection{54} Of this the learned Hesiod did not say so much as a word, though he was writing about the cultivation of the field. But Homer, who lived many ages before, as it seems to me, makes Laertes—soothing the longing he felt for his son—till his field and manure it. And country matters are rich not in grainfields alone and meadows and vineyards and groves, but in gardens too and orchards, and then in the pasturing of flocks, the swarming of bees, the variety of every kind of flower. Nor is it only planting that delights, but grafting too, than which farming has found nothing more ingenious.

\ciceroSection{55} I could go on through very many delights of country life, but I am aware that even what I have said has run on too long. You will forgive me, though—for I have been carried away by my enthusiasm for country matters, and old age is by nature rather talkative, lest I seem to clear it of every fault. It was in this kind of life, then, that Manius Curius, after he had triumphed over the Samnites, the Sabines, and Pyrrhus, spent the final span of his years; and when I contemplate his farmhouse—for it lies not far from mine—I cannot sufficiently admire either the man's own self-restraint or the discipline of those times.

\ciceroSection{56} When the Samnites brought a great weight of gold to Curius as he sat at his hearth, he turned them away; for he said that, as he saw it, the splendid thing was not to possess gold but to give orders to those who possessed it. Could a spirit so great fail to make old age agreeable? But I come back to farmers, so as not to stray from my own kind. In those days senators—that is, old men—lived on the land; for it was while Lucius Quinctius Cincinnatus was at the plow that word reached him he had been made dictator, and it was by that dictator's order that Gaius Servilius Ahala, his master of horse, seized Spurius Maelius, who was reaching for kingship, and put him to death. Curius and the rest of the old men used to be summoned to the Senate from their farmhouses, and that is why the men who summoned them came to be called "travelers." Was the old age of these men, then, a thing to pity—men who took delight in the cultivation of the land? In my judgment, at least, I rather doubt any life can be more blessed, and not only in its duty—since the working of the land is wholesome for the whole human race—but also in the delight I have spoken of, and in the plenty and abundance of all the things that serve men's nourishment and the worship of the gods as well; so that, since some people miss these things, let us now make our peace with pleasure. For a good and diligent master always has his wine cellar stocked, his oil store, his larder too, and the whole farm is rich—overflowing with pig and kid and lamb, with hen, with milk and cheese and honey. The garden, moreover, the farmers themselves call a second flitch of bacon. And to all this, fowling and hunting, in spare hours, add their relish.

\ciceroSection{57} Why should I say more about the green of the meadows, or the rows of trees, or the look of the vineyards and the olive groves? I will cut it short. Nothing can be richer in use or finer in appearance than a well-tended field, and old age, far from holding a man back from its enjoyment, even invites and entices him to it. For where else can that time of life better warm itself, whether by the sun or by the fire, or, in turn, more healthily cool itself in shade and water? Let the young, then, keep their weapons to themselves,

\ciceroSection{58} their horses, their spears, their club and ball, their swimming and their footraces; to us old men, out of their many games, let them leave the dice and the knucklebones—and even those just as we please, since old age can be blessed without them.

\ciceroSection{59} Xenophon's books are useful for many things, and I beg you, read them closely, as you do. How fully he praises the cultivation of the land in that book on the management of an estate, the one titled the \textit{Oeconomicus}! And to make you understand that nothing seemed to him so kingly as the pursuit of farming, in that book Socrates is talking with Critobulus about Cyrus the Younger, king of the Persians, a man outstanding in talent and in the glory of his rule. When Lysander the Spartan, a man of the highest virtue, had come to him at Sardis and brought him gifts from the allies, Cyrus was in every way courteous and kind toward Lysander, and in particular showed him a certain enclosed park, carefully planted. When Lysander marveled at the height of the trees, at the rows laid out in a quincunx, at the soil worked clean of weeds, and at the sweetness of the scents breathing from the flowers, he said that he admired not only the care but also the skill of the man who had measured and arranged it all; and Cyrus replied: "But I measured it all out myself; the rows are mine, the design is mine; and many of those trees were planted by my own hand." Then Lysander, looking at his purple robe, the gleam of his person, and his Persian finery rich with gold and many jewels, said: "Rightly indeed, Cyrus, do they call you blessed, since your good fortune is joined to your virtue!"

\ciceroSection{60} This is the kind of good fortune, then, that old men may enjoy, and age is no bar to our holding fast to our pursuits—both the rest of them and, above all, the cultivation of the land—right up to the very end of old age. We are told that Marcus Valerius Corvinus carried on to his hundredth year, living on his farms and working them when his time of life was already far spent, and that between his first consulship and his sixth there lay forty-six years. And so the span of life our ancestors meant to mark the beginning of old age was matched by the length of his career in office; and his last years were happier than his middle ones, because he had more authority and less labor.

\ciceroSection{61} The crowning glory of old age, indeed, is authority. How great it was in Lucius Caecilius Metellus, how great in Aulus Atilius Calatinus! Of him there is that famous epitaph: "This was the one man whom very many nations agree to have been the first of his people." The whole verse, cut into his tomb, is well known. Rightly, then, was he a man of weight, since the report of all men was at one in his praise. And what a man we lately saw in Publius Crassus, the chief priest, and after him in Marcus Lepidus, who held the same priesthood! Why should I speak of Paulus, or of Africanus, or, as I did before, of Maximus? In them authority resided not only in a pronouncement but even in a nod. Old age, especially when it has been honored, carries authority so great that it outweighs all the pleasures of youth.

\ciceroSection{62} But through this whole discourse, remember that the old age I praise is the one founded on the foundations of youth. From which it follows—as I once said, to the great agreement of all—that the old age which has to defend itself by argument is a wretched one. Gray hairs and wrinkles cannot suddenly snatch authority for themselves; it is the earlier life, honorably spent, that reaps authority as its last fruit.

\ciceroSection{63} For these very things are marks of honor, slight and ordinary as they seem—to be greeted, to be sought out, to have men step aside, to have them rise, to be escorted to and from one's door, to be asked for counsel—observances kept among us, and in other states, more scrupulously the better their morals are. They say Lysander the Spartan, whom I mentioned just now, used to call Sparta the most honorable home for old age; for nowhere is so much rendered to age, nowhere is old age more honored. Indeed, the story has come down to us that once at Athens, during the games, a man well on in years came into the theater, and nowhere in that great gathering was room given him by his own countrymen; but when he came over to the Spartans, who, being there as envoys, were seated in a fixed place, they are all said to have risen and taken the old man in among them to sit.

\ciceroSection{64} And when the whole gathering had given them a thundering round of applause, one of the Spartans said the Athenians know what is right but will not do it. There are many fine things in our college of augurs, but this one we are discussing above all: that as each man is the elder, so he holds the first voice in counsel; and the older augurs are set ahead not only of men superior in rank but even of those who hold command. What pleasures of the body, then, are to be set against the rewards of authority? Those who have used these rewards splendidly seem to me to have played out the drama of life and not, like untrained actors, collapsed in the final act.

\ciceroSection{65} But old men, you will say, are peevish and anxious and quick-tempered and hard to please; and, if we look into it, grasping as well. Yet these are faults of character, not of old age. And still, peevishness and the faults I have named have some excuse—not, to be sure, a just one, but one that seems as if it might be allowed: they think themselves scorned, looked down on, mocked; and besides, in a frail body every slight is hateful. Yet all this is made sweeter by good character and good cultivation, as can be seen both in life and on the stage, from those two brothers in the \textit{Adelphi}. What harshness in the one, what graciousness in the other! So it stands: for just as not every wine, so not every nature turns sour with age. Sternness in old age I approve, but, like everything else, in measure; bitterness, by no means; and what miserliness in an old man is meant to accomplish, I cannot understand.

\ciceroSection{66} For can anything be more absurd than, the less of the road that remains, to seek the more provision for the journey? There remains the fourth charge, the one that seems most of all to distress our time of life and keep it anxious: the approach of death, which surely cannot be far from old age. What a wretched old man, to have failed in so long a life to see that death deserves contempt! For death is either to be utterly disregarded, if it altogether snuffs out the soul, or even to be longed for, if it leads the soul somewhere where it will exist forever. And surely no third possibility can be found.

\ciceroSection{67} Why, then, should I fear, if after death I am to be either not wretched or even blessed? And besides—who is so foolish, however young he may be, that he can be certain he will live until evening? Indeed, that time of life has far more chances of death than ours: the young fall into illness more readily, suffer more grievously, are cured with more difficulty. And so few reach old age; were it not so, life would be lived better and more wisely. For mind and reason and judgment reside in the old, and if there had been no old men, no states at all would ever have existed. But I return to the death that hangs over us. What sort of charge against old age is this, when you see that it is one old age shares with youth?

\ciceroSection{68} I felt, in losing the best of sons, and you, Scipio, in losing brothers marked out for the highest distinction, that death is common to every age. But the young man hopes he will live long, a hope the old man cannot share. He hopes foolishly; for what is more senseless than to treat the uncertain as certain, the false as true? But the old man has nothing even to hope for. Yet he is in a better position than the young man, since what the one only hopes for the other has already attained: the one wants to live long, the other has lived long.

\ciceroSection{69} And yet—good gods—what is there in human nature that lasts long? Grant the fullest span; let us look forward to the age of the king of Tartessus—for there was, as I find it written, a certain Arganthonius of Gades who reigned eighty years and lived a hundred and twenty—still, nothing seems to me long-lasting at all if it has any end. For when that end arrives, then what has passed has flowed away; only this remains, which you have won by virtue and right action. Hours give way, and days, and months, and years; past time never returns, and what is to follow cannot be known. Whatever span of time is granted each man for living, with that he ought to be content.

\ciceroSection{70} For the actor, to win approval, need not play the drama through to the end; he need only prove himself in whatever act he appears. Nor must the wise come all the way to the closing "Applause." A short span of life is long enough for living well and honorably. And if it should run on further, there is no more cause for grief than the farmer feels when, the sweetness of springtime past, summer and autumn have come. For spring, like youth, gives sign and promise of the fruits to come; the seasons that follow are suited to the harvesting and gathering of those fruits.

\ciceroSection{71} And the fruit of old age, as I have often said, is the memory and abundance of goods won before. Now everything that happens in accordance with nature is to be counted among the goods; and what is so much in accordance with nature as for the old to die? The same falls to the young against nature's will and resistance. And so the young seem to me to die as a fire's force is quenched under a flood of water, but the old as a flame goes out, spent of its own accord, with no force applied; and just as apples, when unripe, are pulled from the trees only with difficulty, but when ripe and mellow fall of themselves, so force takes life from the young, ripeness from the old. For me this ripeness is so delightful that, the nearer I come to death, the more I seem to glimpse land, and to be making port at last after a long voyage.

\ciceroSection{72} Old age has no fixed boundary, and a man rightly lives in it as long as he can carry out and uphold the duties of his station and yet hold death in contempt. From this it follows that old age is even more spirited than youth, and braver. This is the meaning of Solon's reply to the tyrant Pisistratus, who asked on what at last he relied to oppose him so boldly; Solon, it is said, answered, "On old age." But the best end of living comes when, the mind whole and the senses sound, nature herself dissolves the work she has joined together. As a ship or a building is most easily taken down by the very builder who put it up, so a man is best dissolved by the very nature that bound him together. Now every fresh joining is hard to break, an old one easily. And so it comes about that the brief remainder of life is neither to be greedily clutched at by the old nor abandoned without cause.

\ciceroSection{73} Pythagoras forbids us to quit the guard-post and station of life without the order of our commander, that is, of the god. There is, too, that epitaph of the wise Solon, in which he says he does not wish his own death to lack the grief and lamentation of his friends. He wants, I suppose, to be dear to his people. But I rather think Ennius does better: Let no one honor me with tears, nor hold my funeral with weeping. He judges that death is not to be mourned when immortality follows upon it.

\ciceroSection{74} Now there may be some sensation in the act of dying, but it lasts only a brief moment, especially for an old man; after death, sensation is either to be longed for or is nothing at all. But this must be rehearsed from youth onward—that we hold death of no account; without that rehearsal no one can be of a tranquil mind. For die we surely must, and it is uncertain whether on this very day. So how can a man who fears death, hanging over him every hour, ever be steady in his soul?

\ciceroSection{75} On this point no very long argument seems needed, when I call to mind not Lucius Brutus, killed in setting his country free; not the two Decii, who spurred their galloping horses to a death freely chosen; not Marcus Atilius, who went out to torture to keep faith pledged to an enemy; not the two Scipios, who were willing to block the Carthaginians' march even with their own bodies; not your grandfather Lucius Paulus, who atoned by his death for the recklessness of his colleague in the disgrace at Cannae; not Marcus Marcellus, whose death not even the cruelest of enemies would let go without the honor of burial—but our own legions, who, as I wrote in my \textit{Origins}, marched off again and again, eager and high in spirit, to a place from which they thought they would never return. What the young hold in contempt—and not only the unschooled among them but even rough countrymen—shall learned old men dread?

\ciceroSection{76} Altogether, as it seems to me at least, a satiety of all pursuits brings on a satiety of life. Boyhood has its own pursuits: do the young, then, long for them? Early youth has its pursuits: does the now-settled age that is called middle ask for them? That age too has its own: not even these are sought in old age. There are, finally, certain pursuits belonging to old age; so, just as the pursuits of the earlier ages fall away, the pursuits of old age fall away as well. And when that comes to pass, a satiety of life brings the ripe time of death.

\ciceroSection{77} For I do not see why I should not dare tell you what I myself feel about death, since I seem to see it the better the nearer I am to it. I judge that your fathers—Publius Scipio, and yours, Gaius Laelius—men most illustrious and most dear to me, are living, and living that life which alone deserves the name of life. For while we are shut within this framework of the body, we are discharging a kind of compulsory duty, a heavy labor; the soul is of heaven, pressed down from its highest dwelling and as it were sunk into the earth, a place opposed to its divine nature and to eternity. But I believe the immortal gods sowed souls into human bodies so that there should be beings to watch over the earth, and who, contemplating the order of the heavens, would imitate it in the measure and constancy of their lives. And it is not reason and argument alone that have driven me to believe this, but also the rank and authority of the greatest philosophers.

\ciceroSection{78} I used to hear that Pythagoras and the Pythagoreans, all but countrymen of ours, since they were once called the Italian philosophers, never doubted that our souls are drawn from the universal divine mind. Moreover, the arguments were laid out for me that Socrates had made on the immortality of souls on the last day of his life — Socrates, judged by the oracle of Apollo to be the wisest of all men. Why labor the point? This is what I have come to believe, this is what I hold: since the speed of the mind is so great, its memory of the past and foresight of the future so great, since there are so many arts, such vast sciences, so many discoveries, the nature that contains all these things cannot be mortal; and since the soul is always in motion and has no beginning of motion, because it moves itself, it will have no end of motion either, because it will never abandon itself; and since the nature of the soul is simple, with nothing mixed into it that is unlike or alien to itself, it cannot be divided, and if it cannot be divided, it cannot perish; and it is powerful evidence that men know most things before they are born that children, even as they learn difficult arts, seize upon countless facts so quickly that they seem not to be taking them in for the first time but to be remembering and recalling them. This, more or less, is Plato's view.

\ciceroSection{79} And in Xenophon the elder Cyrus, as he is dying, says this: "Do not imagine, my dearest sons, that when I have departed from you I will be nowhere or be nothing. For while I was with you, you did not see my soul, yet you understood that it was in this body from the things I did. Believe, then, that it is the same soul still, even though you will see nothing of it.

\ciceroSection{80} Indeed, the honors paid to famous men would not endure after their death if their own souls did nothing to make us hold their memory longer. For my part, I could never be persuaded that souls live while they are in mortal bodies and die when they have left them; nor that the soul becomes mindless once it has escaped a mindless body, but rather that when it is freed from every admixture of the body and has begun to be pure and whole, then it becomes wise. And further, when a man's nature is dissolved by death, it is plain where each of his other parts withdraws to, for everything goes back to where it arose; the soul alone is unseen, both when it is present and when it has departed. And you surely see that nothing is so like death as sleep.

\ciceroSection{81} And yet the souls of sleepers most of all reveal their divinity; for when they are loosed and free, they foresee much that is to come, from which we understand what they will be like once they have wholly slipped the body's chains. Therefore, if these things are so, honor me," he said, "as a god; but if the soul is to perish together with the body, then you, in reverence for the gods who watch over and govern all this beauty, will preserve the memory of me with devotion and without violation." So spoke Cyrus as he died; now let us, if you please, look at our own case.

\ciceroSection{82} No one will ever persuade me, Scipio, that your father Paulus, or your two grandfathers Paulus and Africanus, or the father and uncle of Africanus, or the many distinguished men there is no need to list, attempted such great things — things that bear on the memory of posterity — unless they saw in their minds that posterity belonged to them. Or do you suppose, if I may boast a little about myself, as old men do, that I would have undertaken such great labors by day and night, at home and in the field, if I were to bound my glory by the same limits as my life? Would it not have been far better to pass an idle and quiet old age, free of all toil and struggle? But somehow my soul, lifting itself up, always looked ahead to posterity as though, once it had passed out of life, it would at that point finally begin to live. And surely, were it not so — were souls not immortal — the soul of every best man would not strive hardest of all toward the glory of immortality. And consider that the wisest man dies with the most serene mind,

\ciceroSection{83} the most foolish with the most troubled. Does it not seem to you that the soul which sees more, and sees farther, sees itself setting out for something better, while the one whose vision is duller does not see this? For my part, I am carried away with longing to see your fathers, whom I honored and loved; and not only do I yearn to meet those I knew myself, but also those of whom I have heard and read and have written about myself. And once I have set out toward them, no one will easily drag me back, or boil me young again as they did Pelias. And if some god should grant me to grow young again from this age of mine and cry once more in the cradle, I would refuse it utterly, nor would I want to be called back to the start from the finish line, like a runner whose course is run.

\ciceroSection{84} For what advantage does life hold? What does it not rather hold of toil? But grant that it has its advantages: it certainly has, all the same, either a fullness or a limit. I do not care to lament life, as many men, and learned ones, have often done; nor do I regret having lived, since I have lived in such a way that I do not think I was born in vain; and I leave life as though leaving an inn, not a home — for nature has given us a place to lodge in for a while, not to live in. What a glorious day it will be when I set out for that divine gathering and assembly of souls, and depart from this crowd and corruption! For I shall set out not only to those men of whom I spoke before, but also to my own son Cato, than whom no better man was ever born, none more outstanding in devotion — whose body was burned by me, when it would rightly have been my body burned by him; but his soul, not deserting me but looking back toward me, surely departed for those regions where it saw I myself must come. I seemed to bear that loss of mine bravely — not that I bore it with a calm mind, but I consoled myself with the thought that the parting and separation between us would not be long.

\ciceroSection{85} For these reasons, Scipio — since you said that you and Laelius are accustomed to marvel at it — old age sits lightly on me, and not only is it not a burden, it is even a delight. And if I am wrong in this, in believing that the souls of men are immortal, I am glad to be wrong, and I do not want this error of mine, which gives me such joy, to be wrenched from me while I live; but if, once dead, I shall feel nothing — as certain petty philosophers hold — then I am not afraid that dead philosophers will laugh at this error. And if we are not to be immortal, still it is desirable for a man to be extinguished in his own due time. For nature has a limit to living, as it does to all other things. And old age is the final act of life, as of a play; we should flee its exhaustion, especially once it is joined with satiety. These are the things I had to say about old age. May you both come to it, so that the things you have heard from me you may prove true by your own experience!
